\documentclass[10pt,a4paper]{amsart}
\usepackage{amsmath,amssymb,amsthm,graphicx}
% Retained preamble for the Stacks Project chapter {CH} (shared project preamble).
% The chapter file pinned-original-source.tex only carries \input{{preamble}} on its first line,
% so no preamble text lives inside the chapter itself.  The Stacks Project uses one single shared
% preamble for every chapter: the snapshot stages it three times (sources/stacks-chow/preamble.tex,
% sources/stacks-etale-cohomology/preamble.tex and focused-seeds/AGP/sources/stacks-snapshot-20261011/
% preamble.tex) and those copies are identical in size.  The block below is copied mechanically,
% line for line, from sources/stacks-chow/preamble.tex lines 14-18 and 154-217: the Xy-pic setup and
% the native theorem declarations together with the mathematical macros.  Nothing was added, removed
% or redefined; the wrapper class and the amsmath/amssymb/amsthm/graphicx base packages are supplied
% by the assembler.
% For commutative diagrams we use Xy-pic
% Chapter: stacks-spaces-more-morphisms.tex -- More on Morphisms of Spaces.
\usepackage[all]{xy}

% We use 2cell for 2-commutative diagrams.
\xyoption{2cell}

% Theorem environments.
%
\theoremstyle{plain}
\newtheorem{theorem}[subsection]{Theorem}
\newtheorem{proposition}[subsection]{Proposition}
\newtheorem{lemma}[subsection]{Lemma}

\theoremstyle{definition}
\newtheorem{definition}[subsection]{Definition}
\newtheorem{example}[subsection]{Example}
\newtheorem{exercise}[subsection]{Exercise}
\newtheorem{situation}[subsection]{Situation}

\theoremstyle{remark}
\newtheorem{remark}[subsection]{Remark}
\newtheorem{remarks}[subsection]{Remarks}

\numberwithin{equation}{subsection}

% Macros
%
\def\lim{\mathop{\mathrm{lim}}\nolimits}
\def\colim{\mathop{\mathrm{colim}}\nolimits}
\def\Spec{\mathop{\mathrm{Spec}}}
\def\Hom{\mathop{\mathrm{Hom}}\nolimits}
\def\Ext{\mathop{\mathrm{Ext}}\nolimits}
\def\SheafHom{\mathop{\mathcal{H}\!\mathit{om}}\nolimits}
\def\SheafExt{\mathop{\mathcal{E}\!\mathit{xt}}\nolimits}
\def\Sch{\mathit{Sch}}
\def\Mor{\mathop{\mathrm{Mor}}\nolimits}
\def\Ob{\mathop{\mathrm{Ob}}\nolimits}
\def\Sh{\mathop{\mathit{Sh}}\nolimits}
\def\NL{\mathop{N\!L}\nolimits}
\def\CH{\mathop{\mathrm{CH}}\nolimits}
\def\proetale{{pro\text{-}\acute{e}tale}}
\def\etale{{\acute{e}tale}}
\def\QCoh{\mathit{QCoh}}
\def\Ker{\mathop{\mathrm{Ker}}}
\def\Im{\mathop{\mathrm{Im}}}
\def\Coker{\mathop{\mathrm{Coker}}}
\def\Coim{\mathop{\mathrm{Coim}}}

% Boxtimes
%
\DeclareMathSymbol{\boxtimes}{\mathbin}{AMSa}{"02}

%
% Macros for moduli stacks/spaces
%
\def\QCohstack{\mathcal{QC}\!\mathit{oh}}
\def\Cohstack{\mathcal{C}\!\mathit{oh}}
\def\Spacesstack{\mathcal{S}\!\mathit{paces}}
\def\Quotfunctor{\mathrm{Quot}}
\def\Hilbfunctor{\mathrm{Hilb}}
\def\Curvesstack{\mathcal{C}\!\mathit{urves}}
\def\Polarizedstack{\mathcal{P}\!\mathit{olarized}}
\def\Complexesstack{\mathcal{C}\!\mathit{omplexes}}
% \Pic is the operator that assigns to X its picard group, usage \Pic(X)
% \Picardstack_{X/B} denotes the Picard stack of X over B
% \Picardfunctor_{X/B} denotes the Picard functor of X over B
\def\Pic{\mathop{\mathrm{Pic}}\nolimits}
\def\Picardstack{\mathcal{P}\!\mathit{ic}}
\def\Picardfunctor{\mathrm{Pic}}
\def\Deformationcategory{\mathcal{D}\!\mathit{ef}}
\newenvironment{reference}{}{}
\newenvironment{slogan}{}{}
\newenvironment{history}{}{}
\usepackage{cleveref}
\crefname{originallemma}{original lemma}{original lemmas}
\Crefname{originallemma}{Original lemma}{Original lemmas}
\makeatletter
\newlabel{definition-relative-pseudo-coherence}{{09UI}{source}{}{sourcepointer-1}{}}
\newlabel{definition-relative-pseudo-coherence@cref}{{[originallemma][0][]09UI}{[1][0][]source}{}{}{}}
\makeatother
\crefname{originallemma}{original lemma}{original lemmas}
\Crefname{originallemma}{Original lemma}{Original lemmas}
\makeatletter
\newlabel{lemma-compute-Fourier-Mukai-for-derived-chow}{{0CSK}{source}{}{sourcepointer-2}{}}
\newlabel{lemma-compute-Fourier-Mukai-for-derived-chow@cref}{{[originallemma][0][]0CSK}{[1][0][]source}{}{}{}}
\makeatother
\crefname{originallemma}{original lemma}{original lemmas}
\Crefname{originallemma}{Original lemma}{Original lemmas}
\makeatletter
\newlabel{lemma-derived-chow}{{0CSJ}{source}{}{sourcepointer-3}{}}
\newlabel{lemma-derived-chow@cref}{{[originallemma][0][]0CSJ}{[1][0][]source}{}{}{}}
\makeatother
\crefname{originallemma}{original lemma}{original lemmas}
\Crefname{originallemma}{Original lemma}{Original lemmas}
\makeatletter
\newlabel{more-morphisms-lemma-check-relative-pseudo-coherence-on-charts}{{09UU}{source}{}{sourcepointer-4}{}}
\newlabel{more-morphisms-lemma-check-relative-pseudo-coherence-on-charts@cref}{{[originallemma][0][]09UU}{[1][0][]source}{}{}{}}
\makeatother
\crefname{originallemma}{original lemma}{original lemmas}
\Crefname{originallemma}{Original lemma}{Original lemmas}
\makeatletter
\newlabel{perfect-lemma-pseudo-coherent-on-projective-space}{{0CSG}{source}{}{sourcepointer-5}{}}
\newlabel{perfect-lemma-pseudo-coherent-on-projective-space@cref}{{[originallemma][0][]0CSG}{[1][0][]source}{}{}{}}
\makeatother
\crefname{originallemma}{original lemma}{original lemmas}
\Crefname{originallemma}{Original lemma}{Original lemmas}
\makeatletter
\newlabel{spaces-perfect-lemma-cohomology-base-change}{{08IN}{source}{}{sourcepointer-6}{}}
\newlabel{spaces-perfect-lemma-cohomology-base-change@cref}{{[originallemma][0][]08IN}{[1][0][]source}{}{}{}}
\makeatother
\crefname{originallemma}{original lemma}{original lemmas}
\Crefname{originallemma}{Original lemma}{Original lemmas}
\makeatletter
\newlabel{spaces-perfect-lemma-flat-proper-pseudo-coherent-direct-image-general}{{0CTN}{source}{}{sourcepointer-7}{}}
\newlabel{spaces-perfect-lemma-flat-proper-pseudo-coherent-direct-image-general@cref}{{[originallemma][0][]0CTN}{[1][0][]source}{}{}{}}
\makeatother
\title{Relative pseudo-coherence of a complex on a space over a ring tested against all pseudo-coherent pairings}
\author{The Stacks Project Authors}
\date{stacks-project:spaces-more-morphisms.tex:local-snapshot-20261011; GNU FDL 1.2 or later}
\begin{document}
\maketitle
\noindent\textit{Editorial note.} Statement span is lines 12164-12172 of sources/stacks-spaces-more-morphisms/pinned-original-source.tex; proof span is lines 12174-12227, the complete author proof (52 body lines).\par\medskip
\noindent\textit{Editorial note.} Navigation stubs: 7 cited-but-not-included labels are supplied as non-load-bearing pointers with their pinned tags where the snapshot tag index has them.\par\medskip
\noindent\textit{Editorial note.} extra\_packages is empty on purpose: the retained preamble already loads xy with options and a second declaration clashes in the assembler wrapper.\par\medskip
\noindent\textit{Editorial note.} Text is retained verbatim, including the author's own typos and placeholders.\par\medskip
\noindent\textit{Editorial note.} Counted B side (both complete, spaces more general). Closes L16132 (wave 14) as load-bearing input.\par\medskip
\noindent\textit{Editorial note.} Cross-chapter same-label warning (NOT an occupancy conflict, Chief ruling 2026-10-13 evening): label 'lemma-characterize-pseudo-coherent' also exists in stacks-more-morphisms — a different proposition in a different chapter file.\par\medskip
\noindent\textit{Allowed standard prerequisites (editorial).}\begin{itemize}
\item Explicit prerequisite boundary. Chapter lemma derived-chow (B side, in-source); chapter lemma compute-Fourier-Mukai-for-derived-chow (L16215); Derived Categories of Spaces, lemmas cohomology-base-change and flat-proper-pseudo-coherent-direct-image-general; Derived Categories of Schemes, lemma pseudo-coherent-on-projective-space; More on Morphisms, lemma check-relative-pseudo-coherence-on-charts (L16132).
\end{itemize}\medskip
\ifdefined\hypertarget\hypertarget{sourcepointer-1}{}\fi
\noindent\textit{Source reference.} Original lemma 09UI: sources/stacks-spaces-more-morphisms/pinned-original-source.tex, source lines 10834--10851. Cited by the retained statement or proof but not part of the collected spans; supplied only as a navigation pointer. The argument does not depend on its content.\par\medskip
\ifdefined\hypertarget\hypertarget{sourcepointer-2}{}\fi
\noindent\textit{Source reference.} Original lemma 0CSK: sources/stacks-spaces-more-morphisms/pinned-original-source.tex, source lines 12115--12126. Cited by the retained statement or proof but not part of the collected spans; supplied only as a navigation pointer. The argument does not depend on its content.\par\medskip
\ifdefined\hypertarget\hypertarget{sourcepointer-3}{}\fi
\noindent\textit{Source reference.} Original lemma 0CSJ: sources/stacks-spaces-more-morphisms/pinned-original-source.tex, source lines 12013--12030. Cited by the retained statement or proof but not part of the collected spans; supplied only as a navigation pointer. The argument does not depend on its content.\par\medskip
\ifdefined\hypertarget\hypertarget{sourcepointer-4}{}\fi
\noindent\textit{Source reference.} Original lemma 09UU: sources/stacks-spaces-more-morphisms/pinned-original-source.tex, source lines 12227--12227. Cited by the retained statement or proof but not part of the collected spans; supplied only as a navigation pointer. The argument does not depend on its content.\par\medskip
\ifdefined\hypertarget\hypertarget{sourcepointer-5}{}\fi
\noindent\textit{Source reference.} Original lemma 0CSG: sources/stacks-spaces-more-morphisms/pinned-original-source.tex, source lines 12227--12227. Cited by the retained statement or proof but not part of the collected spans; supplied only as a navigation pointer. The argument does not depend on its content.\par\medskip
\ifdefined\hypertarget\hypertarget{sourcepointer-6}{}\fi
\noindent\textit{Source reference.} Original lemma 08IN: sources/stacks-spaces-more-morphisms/pinned-original-source.tex, source lines 12227--12227. Cited by the retained statement or proof but not part of the collected spans; supplied only as a navigation pointer. The argument does not depend on its content.\par\medskip
\ifdefined\hypertarget\hypertarget{sourcepointer-7}{}\fi
\noindent\textit{Source reference.} Original lemma 0CTN: sources/stacks-spaces-more-morphisms/pinned-original-source.tex, source lines 12227--12227. Cited by the retained statement or proof but not part of the collected spans; supplied only as a navigation pointer. The argument does not depend on its content.\par\medskip
\begin{lemma}
\label{lemma-characterize-pseudo-coherent}
Let $A$ be a ring. Let $X$ be an algebraic space separated and
of finite presentation over $A$. Let $K \in D_\QCoh(\mathcal{O}_X)$.
If $R\Gamma(X, E \otimes^\mathbf{L} K)$ is pseudo-coherent
in $D(A)$ for every pseudo-coherent $E$ in $D(\mathcal{O}_X)$,
then $K$ is pseudo-coherent relative to $A$
(Definition \ref{definition-relative-pseudo-coherence}).
\end{lemma}

\begin{proof}
Assume $K \in D_\QCoh(\mathcal{O}_X)$ and
$R\Gamma(X, E \otimes^\mathbf{L} K)$ is pseudo-coherent
in $D(A)$ for every pseudo-coherent $E$ in $D(\mathcal{O}_X)$.
Let $x \in |X|$. We will show that $K$ is pseudo-coherent relative to $A$
in an \'etale neighbourhood of $x$. This will prove the lemma
by our definition of relative pseudo-coherence.

\medskip\noindent
Choose $n, Z, z, V, E$ as in Lemma \ref{lemma-derived-chow}.
Denote $p : X \times \mathbf{P}^n \to X$ and
$q : X \times \mathbf{P}^n \to \mathbf{P}^n_A$
the projections.
Then for any $i \in \mathbf{Z}$ we have
\begin{align*}
& R\Gamma(\mathbf{P}^n_A,
Rq_*(Lp^*K \otimes^\mathbf{L} E)
\otimes^\mathbf{L}
\mathcal{O}_{\mathbf{P}^n_A}(i)) \\
& =
R\Gamma(X \times \mathbf{P}^n,
Lp^*K \otimes^\mathbf{L} E \otimes^\mathbf{L}
Lq^*\mathcal{O}_{\mathbf{P}^n_A}(i)) \\
& =
R\Gamma(X,
K \otimes^\mathbf{L} Rp_*(E \otimes^\mathbf{L}
Lq^*\mathcal{O}_{\mathbf{P}^n_A}(i)))
\end{align*}
by
Derived Categories of Spaces,
Lemma \ref{spaces-perfect-lemma-cohomology-base-change}.
By
Derived Categories of Spaces, Lemma
\ref{spaces-perfect-lemma-flat-proper-pseudo-coherent-direct-image-general}
the complex $Rp_*(E \otimes^\mathbf{L} Lq^*\mathcal{O}_{\mathbf{P}^n_A}(i))$
is pseudo-coherent on $X$. Hence the assumption tells us the expression
in the displayed formula is a pseudo-coherent object of $D(A)$.
By
Derived Categories of Schemes,
Lemma \ref{perfect-lemma-pseudo-coherent-on-projective-space}
we conclude that $Rq_*(Lp^*K \otimes^\mathbf{L} E)$ is
pseudo-coherent on $\mathbf{P}^n_A$.
By Lemma \ref{lemma-compute-Fourier-Mukai-for-derived-chow}
we have
$$
Rq_*(Lp^*K \otimes^\mathbf{L} E)|_{X \times_A V} =
R(W \to V)_*K|_W
$$
Since $W \to V$ is a closed immersion into an open subscheme of
$\mathbf{P}^n_A$ this means $K|_W$ is pseudo-coherent relative to $A$
for example by
More on Morphisms,
Lemma \ref{more-morphisms-lemma-check-relative-pseudo-coherence-on-charts}.
\end{proof}

\end{document}
